%%%%%%%%%%%%%%%%%%%%%%%%%%%%%%%%%%%%%%%%%%%%%%%%%%%%%%%%%%%%%%%%%%%%%%%%%%%%%%%%%
%	INTRODUCTION
%   Mandated by Project Task 4.2 ("abstract and introduction").
%%%%%%%%%%%%%%%%%%%%%%%%%%%%%%%%%%%%%%%%%%%%%%%%%%%%%%%%%%%%%%%%%%%%%%%%%%%%%%%%%
\lhead{Introduction}
\section{Introduction}
\label{SecIntroduction}

The Traveling Salesman Problem (TSP) asks for the shortest closed tour that
visits every city in a given set exactly once before returning to the
starting city.  Despite its deceptively simple statement, the TSP is
NP-hard~\cite{applegate2006traveling}, and its search space grows
factorially with the number of cities, rendering exhaustive enumeration
infeasible for all but the smallest instances.  As a result, the TSP has
become one of the most widely studied benchmarks in combinatorial
optimisation and a standard testbed for metaheuristic algorithms.

Among metaheuristics, Genetic Algorithms (GAs) are particularly appealing
for the TSP because they maintain a population of candidate solutions and
leverage recombination to explore large regions of the search space
simultaneously~\cite{holland1975adaptation}.  However, applying GAs to
permutation problems introduces a fundamental challenge: conventional
crossover and mutation operators designed for binary or real-valued
representations do not inherently preserve the permutation property.  When
two valid parent tours are recombined, the resulting offspring often contains
duplicate cities and missing cities, producing an \emph{infeasible} solution
whose fitness value is meaningless.

This project investigates how such infeasible solutions arise, designs a
repair mechanism to restore feasibility, and critically analyses the impact
of repair on solution quality, population diversity, and convergence
behaviour.

\subsection{Background}
\label{Intro:Background}

The TSP was first formalised as a mathematical problem in the 1930s and has
since motivated advances in integer programming, branch-and-bound
algorithms, and approximation theory~\cite{applegate2006traveling}.  In the
evolutionary computation community, the TSP serves as a canonical benchmark
because it combines a simple objective (minimise total tour distance) with a
rich constraint structure (each city visited exactly once, forming a
Hamiltonian cycle).

Genetic Algorithms, introduced by Holland~\cite{holland1975adaptation} and
popularised for combinatorial problems by Goldberg~\cite{goldberg1989genetic},
evolve a population of solutions through selection, crossover, and mutation.
For the TSP, solutions are typically encoded as permutations of the city
indices.  Specialised crossover operators --- such as Partially Mapped
Crossover (PMX)~\cite{goldberg1985alleles}, Order Crossover
(OX)~\cite{davis1985applying}, and Cycle Crossover
(CX)~\cite{oliver1987study} --- have been developed to respect the
permutation structure.  Nevertheless, simpler or general-purpose operators
can still produce offspring that violate the ``each city exactly once''
constraint, necessitating some form of constraint handling.

\subsection{Problem Statement}
\label{Intro:ProblemStatement}

When a GA for the TSP employs crossover operators that do not inherently
preserve feasibility (e.g.\ single-point crossover), the resulting offspring
will typically contain duplicate cities at some positions and miss other
cities entirely.  If these infeasible chromosomes are evaluated directly,
their tour-length fitness values are computed over an incomplete or redundant
subset of cities, producing artificially low distances that do not
correspond to any valid tour.  Because selection favours lower fitness
values, infeasible solutions can dominate the population, distorting
selection pressure and steering the evolutionary search away from the
feasible region.

Three broad strategies exist for handling this infeasibility:
\begin{enumerate}[nosep]
    \item \textbf{Feasibility-preserving operators:} design crossover and
          mutation operators that never produce infeasible offspring (e.g.\
          PMX, OX, CX).
    \item \textbf{Penalty functions:} add a penalty term to the fitness of
          infeasible solutions proportional to the degree of constraint
          violation.
    \item \textbf{Repair mechanisms:} detect and correct constraint
          violations in offspring before fitness evaluation.
\end{enumerate}

This project focuses on the third approach --- repair mechanisms --- while
using the other two as baselines for comparative analysis.  The central
research question is: \emph{How does a repair mechanism affect the quality,
diversity, and convergence of a GA applied to the TSP, compared to
alternative constraint-handling strategies?}

\subsection{Contributions}
\label{Intro:Contributions}

The contributions of this report are as follows:
\begin{enumerate}[nosep]
    \item A complete GA implementation for the TSP with permutation
          encoding, multiple selection, crossover, and mutation operators,
          and a configurable repair mechanism supporting three insertion
          strategies (random, positional, nearest-neighbour).
    \item An empirical demonstration of the infeasibility pathology: using
          naive crossover without repair, we show that the GA converges to
          statistically significant but \emph{invalid} fitness values
          ($p < 0.001$, Wilcoxon signed-rank test, 30 independent runs on
          kroA100).
    \item A parameter sensitivity analysis investigating the interaction
          between GA configuration and repair effectiveness.
    \item A comparative evaluation of repair versus alternative
          constraint-handling approaches, with convergence curves, diversity
          tracking, and statistical significance testing.
\end{enumerate}

\subsection{Report Outline}
\label{Intro:Outline}

The remainder of this report is organised as follows.
Section~\ref{SecLiteratureReview} reviews related work on Genetic Algorithms
and constraint handling for the TSP.
Section~\ref{Task1} defines the problem formally and describes the
experimental setup, including the benchmark instance and parameter ranges.
Section~\ref{Task2} presents the GA design and the repair mechanism
implementation.
Section~\ref{Task3} reports the parameter sensitivity analysis,
performance evaluation, and critical analysis of the repair mechanism.
Section~\ref{Task4} provides visualisations of convergence behaviour and
solution quality.
Section~\ref{SecConclusion} concludes the report and discusses directions
for future research.
